\chapter{Massey products}
\label{ch:Massey}

The purpose of this chapter is to provide some general tools,
and to give some specific computations, of Massey products in
$\Ext$.  This material contributes to Table \ref{tab:Massey},
which lists a number of Massey products in $\Ext$ that we need
for various specific purposes.  Most commonly, these Massey
products yield information about Toda brackets via
the Moss Convergence Theorem \ref{thm:Moss}.

We begin with a $\C$-motivic version of a classical theorem of
Adams about symmetric Massey products.


\begin{thm}
\label{thm:Massey-symmetric}
\mbox{}
\begin{enumerate}
\item
If $h_0 x$ is zero, then $\langle h_0, x, h_0 \rangle$ contains 
$\tau h_1 x$.
\item
If $n \geq 1$ and $h_n x$ is zero, then 
$\langle h_n, x, h_n \rangle$ contains $h_{n+1} x$.
\end{enumerate}
\end{thm}

\revv{
\begin{proof}
The element $\Sq^0(h_n) x$ is contained in $\langle h_n, x, h_n \rangle$
\cite{Hirsch55}, where $\Sq^0$ is an algebraic Steenrod operation \cite{May70}.  We compute that $\Sq^0(h_0)$ equals $\tau h_1$
and $\Sq^0(h_n)$ equals $h_{n+1}$ for $n \geq 1$.  These motivic 
computations follow from the analogous classical computations
\cite{Adams60}*{Lemma 2.5.4}.
\end{proof}
}

\revv{
\begin{remark}
The two parts of Theorem \ref{thm:Massey-symmetric} are less different
than they appear.  Because of the specific values of the motivic weights,
we have a factor of $\tau$ in $\Sq^0(h_0)$, while no $\tau$ appears
in $\Sq^0(h_n)$ for $n \geq 1$.
\end{remark}
}



\section{The operator $g$}

The projection map $p:A_* \map A(2)_*$ induces a map
$p_*: \Ext_\C \map \Ext_{A(2)}$.
Because $\Ext_{A(2)}$ is completely known \cite{Isaksen09},
this map is useful for detecting structure in
$\Ext_\C$.  Proposition \ref{prop:g-Massey} provides a tool
for using $p_*$ to compute certain types of Massey products.

\begin{prop}
\label{prop:g-Massey}
Let $x$ be an element of $\Ext_\C$ such that $h_1^4 x = 0$.
Then $p_* \left( \langle h_4, h_1^4, x \rangle \right)$
equals the element $g p_*(x)$ in $\Ext_{A(2)}$.
\end{prop}

\begin{proof}
The idea of the proof is essentially the same as in
\cite{Isaksen19}*{Proposition 3.1}.
The $\Ext_\C$-module $\Ext_{A(2)}$ is a ``Toda module", 
in the sense that Massey products
$\langle x, a, b \rangle$ are defined for all
$x$ in $\Ext_{A(2)}$ and all $a$ and $b$ in $\Ext_\C$ such
that $x \cdot a = 0$ and $ab = 0$.
In particular, the bracket $\langle 1, h_4, h_1^4 \rangle$
is defined in $\Ext_{A(2)}$.
We wish to compute this bracket.

We use the May Convergence Theorem in order to compute the
bracket.  The crossing differentials condition on the theorem
is satisfied because there are no possible differentials that 
could interfere.

The key point is the May differential
$d_4(b_{21}^2) = h_1^4 h_4$.  
This shows that $g$ is contained in $\langle 1, h_4, h_1^4 \rangle$.
Also, the bracket has no indeterminacy by inspection.

Now suppose that $x$ 
is an element of $\Ext_\C$ such that $h_1^4 x = 0$.
Then
\[
p_* \left( \langle h_4, h_1^4, x \rangle \right) =
1 \cdot \langle h_4, h_1^4, x \rangle =
\langle 1, h_4, h_1^4 \rangle \cdot x =
g p_*(x).
\]
\end{proof}

\begin{ex}
\label{ex:g-Massey}
We illustrate the practical usefulness of
Proposition \ref{prop:g-Massey} with a specific
example.  Consider the Massey product
$\langle h_1^3 h_4, h_1, h_2 \rangle$.
The proposition says that
\[
p_* \left( \langle h_1^3 h_4, h_1, h_2 \rangle \right) = h_2 g
\]
in $\Ext_{A(2)}$.  This implies that
$\langle h_1^3 h_4, h_1, h_2 \rangle$ equals $h_2 g$ in
$\Ext_\C$.
\end{ex}

\begin{remark}
The Massey product computation in Example \ref{ex:g-Massey}
is in relatively low dimension, and it can be computed using
other more direct methods.
Table \ref{tab:Massey} lists additional examples, including
some that cannot be determined by more elementary methods.
\end{remark}

\section{The Mahowald operator}
\label{subsctn:Mahowald-operator}

We recall some results from \cite{Isaksen19} about the
Mahowald operator.  The Mahowald operator is defined to be
$M x = \langle x, h_0^3, g_2 \rangle$ for all $x$ such that $h_0^3 x$
equals zero.  As always, one must be cautious about indeterminacy
in $M x$.

There exists a subalgebra $B$ of the $\C$-motivic Steenrod algebra whose
cohomology
$\Ext_B(\M_2, \M_2)$ equals $\M_2[v_3] \otimes_{\M_2} \Ext_{A(2)}$.
The inclusion of $B$ into the $\C$-motivic Steenrod algebra induces a map
$p_*: \Ext_\C \map \Ext_B$.

\begin{prop} \cite{Isaksen19}*{Theorem 1.1}
\label{prop:M-Massey}
The map $p_*: \Ext_\C \map \Ext_B$ takes $M x$ to the product
$(e_0 v_3^2 + h_1^3 v_3^3) p_*(x)$, whenever $M x$ is defined.
\end{prop}

Proposition \ref{prop:M-Massey} is useful in practice for detecting
certain Massey products of the form $\langle x, h_0^3, g_2 \rangle$.
For example, if $x$ is an element of $\Ext_\C$ such that
$h_0^3 x$ equals zero and $e_0 p_*(x)$ is non-zero in $\Ext_{A(2)}$,
then $\langle x, h_0^3, g_2 \rangle$ is non-zero.

\begin{ex}
\label{ex:Mh1}
Proposition \ref{prop:M-Massey} shows that
$\langle h_1, h_0, h_0^2 g_2 \rangle$ is non-zero.  There is
only one non-zero element in the appropriate degree, so we have
identified the Massey product.  We give this element the name
$M h_1$.
\end{ex}

\begin{ex}
\label{ex:M^2h1}
Expanding on Example \ref{ex:Mh1}, 
Proposition \ref{prop:M-Massey} also shows that 
$\langle M h_1, h_0, h_0^2 g_2 \rangle$ is non-zero.
Again, there is only one non-zero element in the appropriate
degree, so we have identified the Massey product.  We give this
element the name $M^2 h_1$.
\end{ex}

\section{Additional computations}

\begin{lemma}
\label{lem:h1^2,h4^2,h1^2,h4^2}
\revdeg{66, 6, 36}
The Massey product $\langle h_1^2, h_4^2, h_1^2, h_4^2 \rangle$
equals $\D_1 h_3^2$.
\end{lemma}

\begin{proof}
Table \ref{tab:Massey} shows that $\D h_2^2$ equals the Massey product
$\langle h_0^2, h_3^2, h_0^2, h_3^2 \rangle$.
Recall the isomorphism between classical $\Ext$ groups
and $\C$-motivic $\Ext$ groups in degrees satisfying $s+f-2w=0$,
as described in Theorem \ref{thm:Chow-zero}.
This shows that 
$\D_1 h_3^2$ equals $\langle h_1^2, h_4^2, h_1^2, h_4^2 \rangle$.
\end{proof}

\rev{
\begin{lemma}
\label{lem:A',h1,h2}
\revdeg{66, 7, 35}
The Massey product $\langle A', h_1, h_2 \rangle$
equals $\tau G_0$.
\end{lemma}

\begin{proof}
Consider the shuffle
\[
A' \langle h_1, h_2, h_1 \rangle =
\langle A', h_1, h_2 \rangle h_1.
\]
Table \ref{tab:Massey} shows that the left side equals
$h_2^2 A'$, which equals $h_1 \cdot \tau G_0$.
This implies that $\langle A', h_1, h_2 \rangle$
contains $\tau G_0$.

The indeterminacy is zero by inspection.
\end{proof}
}

\begin{lemma}
\label{lem:h1^3h4,h1,tgn}
\revdeg{71, 13, 40}
The Massey product $\langle h_1^3 h_4, h_1, \tau g n \rangle$
equals $\tau g^2 n$, with indeterminacy generated by
$M h_0 h_2^2 g$.
\end{lemma}

\begin{proof}
We start by analyzing the indeterminacy.
The product $M c_0 \cdot h_1^3 h_4$ equals
\[
\langle g_2, h_0^3, c_0 \rangle h_1^3 h_4 = 
\langle g_2, h_0^3, h_1^3 h_4 c_0 \rangle =
\langle g_2, h_0^3, h_0 h_2 \cdot h_2 g \rangle =
\langle g_2, h_0^3, h_2 g \rangle h_0 h_2,
\]
which equals $M h_0 h_2^2 g$.
The equalities hold because the indeterminacies are zero, and
the first and last brackets in this computation are given by
Table \ref{tab:Massey}.
This shows that $M h_0 h_2^2 g$ belongs to the indeterminacy.

Table \ref{tab:Massey} shows that
\[
\langle h_2, h_1^3 h_4, h_1 \rangle = 
\langle h_2, h_1, h_1^3 h_4 \rangle
\]
equals $h_2 g$.
Then 
\[
h_2 \langle h_1^3 h_4, h_1, \tau g n \rangle =
\langle h_2, h_1^3 h_4, h_1 \rangle \tau g n = \tau h_2 g^2 n.
\]
This implies that
$\langle h_1^3 h_4, h_1, \tau g n \rangle$
contains either $\tau g^2 n$ or \revv{$\D h_3 g^2$}.
However, the shuffle
\[
h_1 \langle h_1^3 h_4, h_1, \tau g n \rangle =
\langle h_1, h_1^3 h_4, h_1 \rangle \tau g n = 0
\]
eliminates \revv{$\D h_3 g^2$}.
\end{proof}

\begin{lemma}
\label{lem:h3,p',h2}
\revdeg{80, 5, 42}
The Massey product $\langle h_3, p', h_2 \rangle$ equals
$h_0 e_2$, with no indeterminacy.
\end{lemma}

\begin{proof}
We have
\[
\langle h_3, p', h_2 \rangle h_4^2 = 
h_3 \langle p', h_2, h_4^2 \rangle = 
p' \langle h_2, h_4^2, h_3 \rangle.
\]
Table \ref{tab:Massey} shows that the last Massey product
equals $c_2$.
Observe that $p' c_2$ equals $h_0 h_4^2 e_2$.

\revv{Since $h_4^2 \cdot h_6 e_0$ is zero (as usual, we rely
on complete information about classical products in a large range
\cite{Bruner97} \cite{BR20}),}
this shows that $\langle h_3, p', h_2 \rangle$ equals
either $h_0 e_2$ or $h_0 e_2 + h_6 e_0$.
However, shuffle to obtain
\[
\langle h_3, p', h_2 \rangle h_1 = 
h_3 \langle p', h_2, h_1 \rangle,
\]
which must equal zero
\revv{because multiplication by $h_3$ is zero in the appropriate degree.}
Since $h_1(h_0 e_2 + h_6 e_0)$ is non-zero, it cannot equal
$\langle h_3, p', h_2 \rangle$.

The indeterminacy is zero by inspection.
\end{proof}

\begin{remark}
\label{rem:h1^3h4,h1,tgn}
The Massey product of Lemma \ref{lem:h1^3h4,h1,tgn}
cannot be established with Proposition \ref{prop:g-Massey}
because $p_*(\tau g n) = 0$ in $\Ext_{A(2)}$.
\end{remark}

\begin{lemma}
\label{lem:De1+C0,h1^3,h1h4}
\revdeg{82, 12, 45}
The Massey product $\langle \D e_1 + C_0, h_1^3, h_1 h_4 \rangle$
equals $(\D e_1 + C_0) g$, with no indeterminacy.
\end{lemma}

\begin{proof}
Consider the Massey product
$\langle \tau (\D e_1 + C_0), h_1^4, h_4 \rangle$.
By inspection, this Massey product has no indeterminacy. 
Therefore,
\[
\langle \tau (\D e_1 + C_0), h_1^4, h_4 \rangle =
(\D e_1 + C_0) \langle \tau, h_1^4, h_4 \rangle.
\]
Table \ref{tab:Massey} shows that the latter bracket equals $\tau g$,
so the expression equals $\tau (\D e_1 + C_0) g$.

On the other hand, it also equals
$\tau \langle \D e_1 + C_0, h_1^3, h_1 h_4 \rangle$.
Therefore, the bracket
$\langle \D e_1 + C_0, h_1^3, h_1 h_4 \rangle$ must contain
$(\D e_1 + C_0) g$.  Finally, the indeterminacy can be computed
by inspection.
\end{proof}


\begin{lemma}
\label{lem:t^3gG0,h0h2,h2}
\revdeg{93, 13, 49}
The Massey product $\langle \tau^3 g G_0, h_0 h_2, h_2 \rangle$
has indeterminacy generated by $\tau M^2 h_2$, and it either
contains zero or $\tau e_0 x_{76,9}$.  In particular, it does
not contain any linear combination of $\D^2 h_1 g_2$
with other elements.
\end{lemma}

\begin{proof}
The indeterminacy can be computed by inspection.

The only possible elements in
the Massey product $\langle \tau^2 g G_0, h_0 h_2, h_2 \rangle$
are linear combinations of $e_0 x_{76,9}$ and $M^2 h_2$.
The inclusion
\[
\tau \langle \tau^2 g G_0, h_0 h_2, h_2 \rangle \subseteq
\langle \tau^3 g G_0, h_0 h_2, h_2 \rangle
\]
gives the desired result.
\end{proof}


